\subsection{Bias-Variance and Approximation-Estimation Tradeoffs}

Learning quality depends both on the richness of the model class and on the variability introduced by fitting that class from a finite sample.

\begin{proposition} \label{prop:sl-bias-variance-regression} \hlt{Bias-Variance Decomposition}\\
Let $Y=f_0(X)+\varepsilon$ with $\mathbb{E}[\varepsilon\mid X=x]=0$ and $\operatorname{Var}(\varepsilon\mid X=x)=\sigma^2(x)$.\\
Let random $S$ produce $\widehat{f}_S$, set $\bar{f}(x)\defeq\mathbb{E}_S\widehat{f}_S(x)$, and let $Y_x\perp S$ be a fresh response at $x$.\\
Then $\mathbb{E}_{S,Y_x}[(Y_x-\widehat{f}_S(x))^2]=\sigma^2(x)+(f_0(x)-\bar{f}(x))^2+\operatorname{Var}_S(\widehat{f}_S(x))$.
\begin{proof} Write $Y_x-\widehat{f}_S(x)=\varepsilon_x+f_0(x)-\widehat{f}_S(x)$, where $\varepsilon_x\defeq Y_x-f_0(x)$.\\
Since $\mathbb{E}[\varepsilon_x]=0$ and $Y_x$ is conditionally independent of $S$, $\mathbb{E}_{S,Y_x}[\varepsilon_x\{f_0(x)-\widehat{f}_S(x)\}]=0$.\\
Thus $\mathbb{E}_{S,Y_x}[(Y_x-\widehat{f}_S(x))^2]=\sigma^2(x)+\mathbb{E}_S[(f_0(x)-\widehat{f}_S(x))^2]$.\\
Finally, $\mathbb{E}_S[(f_0(x)-\widehat{f}_S(x))^2]=(f_0(x)-\bar{f}(x))^2+\mathbb{E}_S[(\widehat{f}_S(x)-\bar{f}(x))^2]$.\\
\hlt{Structural Commentary}: Squared prediction error equals noise plus fitted-function error.\\
\hlt{Key Insight}: Center $\widehat{f}_S(x)$ at $\bar{f}(x)$; the remaining cross term has mean zero.
\end{proof}
\end{proposition}

\begin{figure}[H]
  \centering
  \begin{tikzpicture}[
  x=0.72cm,
  y=0.72cm,
  >=Stealth,
  line cap=round,
  line join=round,
  every node/.style={font=\scriptsize},
  axis/.style={->, gray!70, line width=0.45pt},
  train/.style={blue!70!black, line width=1.0pt},
  test/.style={red!75!black, line width=1.0pt},
  guide/.style={gray!55, dashed, line width=0.45pt}
]
  \path[use as bounding box] (-0.55,-0.58) rectangle (7.35,3.95);
  \draw[axis] (0,0) -- (6.7,0) node[black!75,right] {complexity};
  \draw[axis] (0,0) -- (0,3.6) node[black!75,above] {error};

  \draw[train] plot[smooth] coordinates {
    (0.35,3.05) (1.10,2.20) (2.00,1.45) (3.00,0.92)
    (4.15,0.55) (5.35,0.33) (6.20,0.22)
  };

  \draw[test] plot[smooth] coordinates {
    (0.35,3.20) (1.10,2.25) (2.00,1.50) (2.85,1.05)
    (3.55,0.90) (4.25,0.98) (5.10,1.35) (6.20,2.20)
  };

  \draw[guide] (3.55,0) -- (3.55,0.90);
  \node[below=2pt, text=black!80] at (3.55,0) {$c^\star$};
\end{tikzpicture}

  \caption{Bias--variance tradeoff: \textcolor{blue!70!black}{blue}: training error, \textcolor{red!75!black}{red}: test error; dashed line marks test minimum.}
  \label{fig:sl-bias-variance-model-complexity}
\end{figure}

\begin{method} \label{meth:sl-nearest-neighbor-local-averaging} \hlt{Nearest-Neighbor Local Averaging}\\
Let $\mathcal{T}\defeq\{(x_i,y_i):1\leq i\leq n\}$ be training data in a metric space $(\mathcal{X},d)$.\\
Let $N_k(x)\subseteq\{1,\ldots,n\}$ be the indices of the $k$ smallest distances $d(x_i,x)$.\\
For numerical responses, define the local average $\widehat{f}_k(x)\defeq \frac{1}{k}\sum\limits_{i\in N_k(x)}y_i$.\\
For class labels $g_i\in\mathcal{G}$, define $\widehat{p}_g(x)\defeq \frac{1}{k}\sum\limits_{i\in N_k(x)}\mathbf{1}_{\{g_i=g\}}$ and classify by $\widehat{g}_k(x)\in\operatorname*{argmax}\limits_{g\in\mathcal{G}}\widehat{p}_g(x)$.\\
Thus $k$-nearest neighbours replace conditioning at $X=x$ by averaging over $N_k(x)$.
\end{method}

\begin{remark} \label{rem:sl-nearest-neighbors-dimensionality} \hlt{Locality in High Dimension}\\
The parameter $k$ trades local bias against variance: small $k$ reacts strongly to $S$, while large $k$ averages more but may wash out local structure.
In $[0,1]^p$, a subcube with volume fraction $r$ has side length $e_p(r)=r^{1/p}$, so $e_{10}(0.01)\approx0.63$ and $e_{10}(0.10)\approx0.80$.
For $N$ uniform points in the unit ball, the median nearest-neighbour radius is $d(p,N)=(1-2^{-1/N})^{1/p}$, about $0.52$ for $(p,N)=(10,500)$.
Thus high-dimensional neighbourhoods large enough to contain data may not be local; useful local methods need structure or an adapted metric.
\end{remark}

\begin{definition} \label{def:sl-approximation-estimation-tradeoff} \hlt{Approximation and Estimation Errors}\\
Let $R^*\defeq\inf_fR(f)$ and $R_{\mathcal{F}}^*\defeq\inf_{f\in\mathcal{F}}R(f)$.
\begin{enumerate}[label=(\roman*)]
\setlength{\itemsep}{0pt}
\item \hlt{Approximation error}: $\mathcal{A}(\mathcal{F})\defeq R_{\mathcal{F}}^*-R^*$.
\item \hlt{Estimation error}: For a data-dependent $\widehat{f}\in\mathcal{F}$, $\mathcal{E}(\widehat{f};\mathcal{F})\defeq R(\widehat{f})-R_{\mathcal{F}}^*$.
\item \hlt{Excess risk}: $R(\widehat{f})-R^*=\mathcal{A}(\mathcal{F})+\mathcal{E}(\widehat{f};\mathcal{F})$.
\end{enumerate}
Larger $\mathcal{F}$ usually lowers $\mathcal{A}(\mathcal{F})$ and raises estimation difficulty.
\end{definition}

\begin{remark} \label{rem:sl-model-complexity-tradeoff} \hlt{Model Complexity}\\
As $c(\mathcal{F})$ grows, $\widehat{R}_S(\widehat{f}_{\mathcal{F}})$ usually decreases, while $R(\widehat{f}_{\mathcal{F}})-R^*=\mathcal{A}(\mathcal{F})+\mathcal{E}(\widehat{f}_{\mathcal{F}};\mathcal{F})$ trades smaller approximation error for larger estimation error.
For zero-one loss, the relevant failure is crossing the wrong side of the decision boundary, not an additive squared-error bias--variance split.
\end{remark}
